% ═══════════════════════════════════════════════════════════════════════════════
\section{Übersetzungsregeln und Operatoren}
\label{sec:operatoren}
% ═══════════════════════════════════════════════════════════════════════════════

Dieses Kapitel definiert die vollständigen Übersetzungsregeln von den
Quellsprachen Python, Java und C\# in den IQT sowie vom IQT nach ANSI SQL.
Die Regeln werden als Transformationstabellen und formale Schreibweisen
dargestellt.

% ─────────────────────────────────────────────────────────────────────────────
\subsection{Operator 1: Filterung (WHERE / HAVING)}
\label{subsec:op_filter}
% ─────────────────────────────────────────────────────────────────────────────

\begin{sidewaystable}
\centering
\caption{Übersetzungstabelle: Filteroperatoren}
\label{tab:filter_ops}
\renewcommand{\arraystretch}{1.4}
\begin{tabularx}{\textwidth}{lXXX}
\toprule
\textbf{IQT-Prädikat} & \textbf{Python} & \textbf{Java} & \textbf{C\#} \\
\midrule
$\mathtt{EQ}(a, v)$
  & \texttt{Model.col == v}
  & \texttt{cb.equal(r.get("col"),v)}
  & \texttt{x.Col == v} \\
$\mathtt{NEQ}(a, v)$
  & \texttt{Model.col != v}
  & \texttt{cb.notEqual(...)}
  & \texttt{x.Col != v} \\
$\mathtt{GT}(a, v)$
  & \texttt{Model.col > v}
  & \texttt{cb.gt(...)}
  & \texttt{x.Col > v} \\
$\mathtt{GTE}(a, v)$
  & \texttt{Model.col >= v}
  & \texttt{cb.ge(...)}
  & \texttt{x.Col >= v} \\
$\mathtt{IN}(a, list)$
  & \texttt{Model.col.in\_(list)}
  & \texttt{root.get("col").in(list)}
  & \texttt{list.Contains(x.Col)} \\
$\mathtt{BETWEEN}$
  & \texttt{col.between(a,b)}
  & \texttt{cb.between(r.get("col"),a,b)}
  & \texttt{x.Col >= a \&\& x.Col <= b} \\
$\mathtt{LIKE}(a, p)$
  & \texttt{Model.col.like(p)}
  & \texttt{cb.like(r.get("col"),p)}
  & \texttt{EF.Functions.Like(x.Col,p)} \\
$\mathtt{IS\_NULL}$
  & \texttt{Model.col == None}
  & \texttt{cb.isNull(r.get("col"))}
  & \texttt{x.Col == null} \\
$\mathtt{AND}$
  & \texttt{\&} oder \texttt{and\_()}
  & \texttt{cb.and(...)}
  & \texttt{\&\&} \\
$\mathtt{OR}$
  & \texttt{|} oder \texttt{or\_()}
  & \texttt{cb.or(...)}
  & \texttt{||} \\
\bottomrule
\end{tabularx}
\end{sidewaystable}

\textbf{SQL-Generierungsregel für FILTER:}
\begin{align*}
\mathrm{SQL}(\mathtt{FILTER}(S, p, \mathtt{WHERE}))
  &= \mathrm{SQL}(S) + \texttt{ WHERE } + \mathrm{SQL}(p) \\
\mathrm{SQL}(\mathtt{FILTER}(S, p, \mathtt{HAVING}))
  &= \mathrm{SQL}(S) + \texttt{ HAVING } + \mathrm{SQL}(p)
\end{align*}

% ─────────────────────────────────────────────────────────────────────────────
\subsection{Operator 2: Projektion (SELECT)}
\label{subsec:op_project}
% ─────────────────────────────────────────────────────────────────────────────

\begin{sidewaystable}
\centering
\caption{Übersetzungstabelle: Projektionsoperatoren}
\label{tab:project_ops}
\renewcommand{\arraystretch}{1.4}
\begin{tabularx}{\textwidth}{lXXX}
\toprule
\textbf{Operation} & \textbf{Python} & \textbf{Java} & \textbf{C\#} \\
\midrule
Alle Spalten
  & \texttt{select(Model)}
  & \texttt{cq.select(root)}
  & \texttt{.Select(x => x)} \\
Einzelne Spalten
  & \texttt{select(M.a, M.b)}
  & \texttt{cb.construct(DTO, r.get("a"), r.get("b"))}
  & \texttt{.Select(x => new\{x.A,x.B\})} \\
Ausdruck
  & \texttt{select(func.upper(M.name))}
  & \texttt{cb.upper(root.get("name"))}
  & \texttt{.Select(x => x.Name.ToUpper())} \\
Distinct
  & \texttt{select(M.col).distinct()}
  & \texttt{cq.distinct(true)}
  & \texttt{.Select(x => x.Col).Distinct()} \\
\bottomrule
\end{tabularx}
\end{sidewaystable}

% ─────────────────────────────────────────────────────────────────────────────
\subsection{Operator 3: Sortierung (ORDER BY)}
\label{subsec:op_sort}
% ─────────────────────────────────────────────────────────────────────────────

\begin{sidewaystable}
\centering
\caption{Übersetzungstabelle: Sortieroperatoren}
\label{tab:sort_ops}
\renewcommand{\arraystretch}{1.4}
\begin{tabularx}{\textwidth}{lXXX}
\toprule
\textbf{IQT-Operation} & \textbf{Python} & \textbf{Java} & \textbf{C\#} \\
\midrule
ASC
  & \texttt{order\_by(M.col)}
  & \texttt{cb.asc(r.get("col"))}
  & \texttt{OrderBy(x => x.Col)} \\
DESC
  & \texttt{order\_by(desc(M.col))}
  & \texttt{cb.desc(r.get("col"))}
  & \texttt{OrderByDescending(x => x.Col)} \\
Multi-Sort
  & \texttt{order\_by(M.a, desc(M.b))}
  & \texttt{.orderBy(cb.asc(...), cb.desc(...))}
  & \texttt{.OrderBy(x => x.A).ThenByDescending(x => x.B)} \\
NULLS FIRST
  & \texttt{nullsfirst(M.col)}
  & \texttt{cb.asc(r.get("col")).nullsFirst()}
  & \texttt{EF nicht nativ} \\
\bottomrule
\end{tabularx}
\end{sidewaystable}

\textbf{SQL-Generierungsregel:}
\[
\mathrm{SQL}(\mathtt{SORT}(S, [s_1, \ldots, s_n])) =
  \mathrm{SQL}(S) + \texttt{ ORDER BY } +
  \mathrm{SQL}(s_1) + \texttt{, } + \ldots + \texttt{, } + \mathrm{SQL}(s_n)
\]
\[
\mathrm{SQL}(\mathtt{SortItem}(e, \mathtt{ASC}, \mathtt{DEFAULT})) =
  \mathrm{SQL}(e) + \texttt{ ASC}
\]

% ─────────────────────────────────────────────────────────────────────────────
\subsection{Operator 4: Begrenzung (LIMIT / OFFSET)}
\label{subsec:op_limit}
% ─────────────────────────────────────────────────────────────────────────────

\begin{sidewaystable}
\centering
\caption{Übersetzungstabelle: LIMIT / OFFSET je Datenbankdialekt}
\label{tab:limit_ops}
\renewcommand{\arraystretch}{1.4}
\begin{tabularx}{\textwidth}{lXX}
\toprule
\textbf{Sprache / Dialekt} & \textbf{Quellausdruck} & \textbf{SQL-Ausgabe} \\
\midrule
Python (SA)
  & \texttt{.limit(10).offset(20)}
  & \texttt{LIMIT 10 OFFSET 20} \\
Java (JPA)
  & \texttt{.setMaxResults(10).setFirstResult(20)}
  & \texttt{LIMIT 10 OFFSET 20} \\
C\# (LINQ)
  & \texttt{.Take(10).Skip(20)}
  & \texttt{LIMIT 10 OFFSET 20} \\
\midrule
\multicolumn{3}{l}{\textit{Dialektspezifische Ausgabe (Erweiterungsmodus):}} \\
MSSQL
  & -- (wie oben)
  & \texttt{OFFSET 20 ROWS FETCH NEXT 10 ROWS ONLY} \\
Oracle $<$12c
  & -- (wie oben)
  & \texttt{WHERE ROWNUM BETWEEN 21 AND 30} (umgeschrieben) \\
\bottomrule
\end{tabularx}
\end{sidewaystable}

% ─────────────────────────────────────────────────────────────────────────────
\subsection{Operator 5: Aggregation (GROUP BY)}
\label{subsec:op_agg}
% ─────────────────────────────────────────────────────────────────────────────

\begin{sidewaystable}
\centering
\caption{Übersetzungstabelle: Aggregationsfunktionen}
\label{tab:agg_ops}
\renewcommand{\arraystretch}{1.4}
\begin{tabularx}{\textwidth}{lXXX}
\toprule
\textbf{IQT-Funktion} & \textbf{Python} & \textbf{Java} & \textbf{C\#} \\
\midrule
\texttt{COUNT(*)}
  & \texttt{func.count('*')}
  & \texttt{cb.count(root)}
  & \texttt{.Count()} \\
\texttt{COUNT(col)}
  & \texttt{func.count(M.col)}
  & \texttt{cb.count(r.get("col"))}
  & \texttt{.Count(x => x.Col != null)} \\
\texttt{SUM(col)}
  & \texttt{func.sum(M.col)}
  & \texttt{cb.sum(r.get("col"))}
  & \texttt{.Sum(x => x.Col)} \\
\texttt{AVG(col)}
  & \texttt{func.avg(M.col)}
  & \texttt{cb.avg(r.get("col"))}
  & \texttt{.Average(x => x.Col)} \\
\texttt{MIN(col)}
  & \texttt{func.min(M.col)}
  & \texttt{cb.min(r.get("col"))}
  & \texttt{.Min(x => x.Col)} \\
\texttt{MAX(col)}
  & \texttt{func.max(M.col)}
  & \texttt{cb.max(r.get("col"))}
  & \texttt{.Max(x => x.Col)} \\
GROUP BY
  & \texttt{group\_by(M.dep)}
  & \texttt{cq.groupBy(r.get("dep"))}
  & \texttt{.GroupBy(x => x.Dep)} \\
HAVING
  & \texttt{having(func.count('*') > 5)}
  & \texttt{cq.having(cb.gt(cb.count(r),5))}
  & \texttt{.Where(g => g.Count() > 5)} \\
\bottomrule
\end{tabularx}
\end{sidewaystable}

% ─────────────────────────────────────────────────────────────────────────────
\subsection{Operator 6: Joins}
\label{subsec:op_join}
% ─────────────────────────────────────────────────────────────────────────────

\begin{sidewaystable}
\centering
\caption{Übersetzungstabelle: Join-Typen}
\label{tab:join_ops}
\renewcommand{\arraystretch}{1.4}
\begin{tabularx}{\textwidth}{lXXX}
\toprule
\textbf{Join-Typ} & \textbf{Python} & \textbf{Java} & \textbf{C\#} \\
\midrule
INNER JOIN
  & \texttt{join(Order, User.id == Order.user\_id)}
  & \texttt{root.join("orders")}
  & \texttt{join in context.Orders on u.Id equals o.UserId} \\
LEFT OUTER
  & \texttt{outerjoin(Order, ...)}
  & \texttt{root.join(..., JoinType.LEFT)}
  & \texttt{join o in ... into grp from o in grp.DefaultIfEmpty()} \\
RIGHT OUTER
  & \texttt{(SA kein nativer Support)}
  & \texttt{JoinType.RIGHT}
  & \texttt{(symmetrisch zu LEFT)} \\
CROSS JOIN
  & \texttt{select\_from(A, B)}
  & \texttt{cq.from(A.class); cq.from(B.class)}
  & \texttt{from a in As from b in Bs select ...} \\
\bottomrule
\end{tabularx}
\end{sidewaystable}

\begin{regel}[JOIN-Übersetzung]
\label{regel:join}
Für einen \texttt{INNER JOIN} zwischen Relation $L$ (Alias $t_1$) und $R$ (Alias $t_2$)
mit Prädikat $p$ ergibt sich:
\[
\mathrm{SQL}(\mathtt{JOIN}(L, R, \mathtt{INNER}, p)) =
  \mathrm{SQL}(L) + \texttt{ INNER JOIN } + \mathrm{SQL}(R) +
  \texttt{ ON } + \mathrm{SQL}(p)
\]
\end{regel}

% ─────────────────────────────────────────────────────────────────────────────
\subsection{Operator 7: Fensterfunktionen (OVER)}
\label{subsec:op_window}
% ─────────────────────────────────────────────────────────────────────────────

Fensterfunktionen (Window Functions, ANSI SQL:2003+) werden in UQTL ab
Konformitätsstufe CL-3 unterstützt.

\begin{sidewaystable}
\centering
\caption{Übersetzungstabelle: Fensterfunktionen}
\label{tab:window_ops}
\renewcommand{\arraystretch}{1.4}
\begin{tabularx}{\textwidth}{lXXX}
\toprule
\textbf{SQL-Funktion} & \textbf{Python} & \textbf{Java} & \textbf{C\#} \\
\midrule
\texttt{ROW\_NUMBER()}
  & \texttt{func.row\_number().over(...)}
  & \texttt{cb.rowNumber().over(...)}
  & \texttt{EF.Functions.RowNumber(...)} \\
\texttt{RANK()}
  & \texttt{func.rank().over(...)}
  & \texttt{cb.rank().over(...)}
  & \texttt{EF.Functions.Rank(...)} \\
\texttt{DENSE\_RANK()}
  & \texttt{func.dense\_rank().over(...)}
  & \texttt{cb.denseRank().over(...)}
  & \texttt{EF.Functions.DenseRank(...)} \\
\texttt{LAG(col, n)}
  & \texttt{func.lag(M.col, n).over(...)}
  & \texttt{cb.lag(r.get("col"), n)}
  & \texttt{EF.Functions.Lag(x.Col, n)} \\
\texttt{LEAD(col, n)}
  & \texttt{func.lead(M.col, n).over(...)}
  & \texttt{cb.lead(r.get("col"), n)}
  & \texttt{EF.Functions.Lead(x.Col, n)} \\
\texttt{SUM() OVER}
  & \texttt{func.sum(M.col).over(...)}
  & \texttt{cb.sum(r.get("col")).over(...)}
  & \texttt{.Sum(x.Col).Over(...)} \\
\bottomrule
\end{tabularx}
\end{sidewaystable}

% ─────────────────────────────────────────────────────────────────────────────
\subsection{Operator 8: Unterabfragen und CTEs}
\label{subsec:op_subquery}
% ─────────────────────────────────────────────────────────────────────────────

\begin{sidewaystable}
\centering
\caption{Übersetzungstabelle: Unterabfragen}
\label{tab:subquery_ops}
\renewcommand{\arraystretch}{1.4}
\begin{tabularx}{\textwidth}{lXXX}
\toprule
\textbf{Typ} & \textbf{Python} & \textbf{Java} & \textbf{C\#} \\
\midrule
Scalar Subquery
  & \texttt{subquery()}
  & \texttt{cb.subquery(Long.class)}
  & \texttt{context.Baz.Where(...).Select(x => x.Val).FirstOrDefault()} \\
EXISTS
  & \texttt{exists(subq)}
  & \texttt{cb.exists(subq)}
  & \texttt{.Any(cond)} \\
IN (Subquery)
  & \texttt{M.col.in\_(subq)}
  & \texttt{root.get("col").in(subq)}
  & \texttt{list.Contains(x.Col)} (mit AsQueryable) \\
CTE (WITH)
  & \texttt{cte = subq.cte("name")}
  & \texttt{(JPA kein nativer CTE)}
  & \texttt{(EF Core $\geq$7: FromSqlRaw)} \\
\bottomrule
\end{tabularx}
\end{sidewaystable}

% ─────────────────────────────────────────────────────────────────────────────
\subsection{Operator 9: Mengenoperationen}
\label{subsec:op_set}
% ─────────────────────────────────────────────────────────────────────────────

\begin{table}[H]
\centering
\caption{Übersetzungstabelle: Mengenoperationen}
\label{tab:set_ops}
\renewcommand{\arraystretch}{1.4}
\begin{tabularx}{\textwidth}{lXXX}
\toprule
\textbf{IQT} & \textbf{Python} & \textbf{Java} & \textbf{C\#} \\
\midrule
UNION
  & \texttt{q1.union(q2)}
  & \texttt{(via CriteriaQuery)} 
  & \texttt{q1.Union(q2)} \\
UNION ALL
  & \texttt{q1.union\_all(q2)}
  & \texttt{(JPA: JPQL UNION)} 
  & \texttt{q1.Concat(q2)} \\
INTERSECT
  & \texttt{q1.intersect(q2)}
  & \texttt{(selten nativ)}
  & \texttt{q1.Intersect(q2)} \\
EXCEPT
  & \texttt{q1.except\_(q2)}
  & \texttt{(selten nativ)}
  & \texttt{q1.Except(q2)} \\
\bottomrule
\end{tabularx}
\end{table}
